\clearpage
\fancyhead[L]{\small JUNO. ENTDECKUNG.}
\hypertarget{native-258-lower}{}
\begin{center}
EINIGE NACHRICHTEN ÜBER DEN NEUEN PLANETEN.

\rule{0.75\linewidth}{0.4pt}

Braunschweigisches Magazin. Stück 40. 1804 October 6.

\rule{0.75\linewidth}{0.4pt}
\end{center}

Die vor kurzen zu \emph{Lilienthal} gemachte Entdeckung eines neuen beweglichen Sterns, von welchem bereits einige Zeitungsnachrichten ins Publicum gekommen sind, hat die öffentliche Neugierde so sehr rege gemacht, dass wahrscheinlich einige bestimmtere Nachrichten über diese merkwürdige Entdeckung den Freunden der Fortschritte der Wissenschaften nicht unwillkommen sein werden; zumal da sich
\SourceLeaf{259}
jetzt, nachdem der Lauf des Sterns seit drei Wochen beobachtet ist, schon ziemlich zuverlässige Resultate über seine wahre Bahn ziehen lassen, die sonst den meisten Lesern dieser Blätter erst weit später zu Gesicht kommen würden.

Herr \textsc{Harding}, Mitarbeiter des berühmten Justizraths \textsc{Schröter} zu Lilienthal arbeitete bereits seit einem halben Jahre an einem astronomischen Atlas eigener Art, welcher nur diejenigen Theile des Fixsternhimmels umfassen soll, in denen die beiden neuentdeckten Planete \emph{Ceres} und \emph{Pallas} sichtbar sein können. Die Aufsuchung und Beobachtung dieser kleinen Weltkörper erfordert durchaus sehr detaillirte Sternkarten, die wenigstens alle Sterne bis zur achten Grösse herab vollständig enthalten, und alle bisher vorhandenen Himmelskarten, ohne Ausnahme, sind dazu bei weitem nicht hinlänglich. Herrn \textsc{Harding}'s Unternehmen ist daher sehr verdienstlich, und von der grossen Sorgfalt und dem unermüdeten Fleisse, womit er dasselbe ausführt, sind einige schon von ihm gelieferte, zu Herrn von \textsc{Zach}'s Monatlicher Correspondenz gehörige vortreffliche Specialkarten Bürgen. Hr. \textsc{Harding} legt dabei hauptsächlich die Histoire céleste française von \textsc{Lalande} zum Grunde, ein unschätzbares Werk, das die zu Paris auf der Sternwarte der Militairschule,
 hauptsächlich von \textsc{Lalande}'s Neffen \textsc{Michel Lefrançais Lalande} binnen 10 Jahren gemachten Beobachtungen von 50000 Fixsternen enthält. Alle diese Sterne, so viele deren in die von Hr. \textsc{Harding} bearbeiteten Zonen fallen, trägt er in seine Karten ein; aber bei jedem einzelnen sieht er vorher selbst am Himmel nach, um sich seines wirklichen Vorhandenseins zu versichern, und die etwaigen bei einer so ungeheuren Arbeit unvermeidlichen Irrthümer, Beobachtungs-, Schreib- und Druckfehler zu verbessern. Aber damit begnügt sich Hr. \textsc{Harding} noch nicht. Überall, wo er noch kleinere Fixsterne, bis zur achten Grösse wenigstens, antrifft, die in der Histoire céleste nicht mit vorkommen, und solcher gibts noch eine grosse Menge, bestimmt er sie selbst, und trägt sie in seinen Karten nach.

Gerade bei einer solchen Durchmusterung einer Himmelsgegend war es, dass Hr. \textsc{Harding} am 1. September im Sternbilde der Fische einen Stern achter Grösse antraf, der nicht in der Histoire céleste steht, und den er daher, ohne Arges daraus zu haben, weil der Fall unzähligemale vorkommt, als einen Fixstern, vorläufig blos nach dem Augenmasse, in seine Karten eintrug. Allein drei Tage später, als er von neuem nachsah, stand auf jenem Platze kein Stern mehr; dagegen in der Nähe, etwa einen halben Grad südwestwärts davon, wo er sich nicht erinnerte, am 1. September einen Stern gesehen zu haben, zeigte sich ein ganz ähnlicher. Dies erregte grossen Verdacht, berechtigte aber noch zu keinem gewissen Schlusse, weil es vielleicht auch nur die Folge eines bei kleinen Fixsternen sehr gewöhnlichen Lichtwechsels, oder eines Irrthums sein konnte. Endlich am 5. September waren die beiden vorigen Plätze leer, und dagegen ein in proportionirter Distanz nach Südwesten gelegener besetzt. Herr \textsc{Harding} konnte nun an der Identität der an den drei Abenden wahrgenommenen Sterne nicht mehr zweifeln. Er beobachtete also den entdeckten beweglichen Stern förmlich, setzte an den folgenden Abenden seine Observationen, die eine ganz regelmässig rückläufige und zugleich nach Süden gehende Bewegung bestätigten, fort, und gab auswärtigen Astronomen von seiner wichtigen Entdeckung Nachricht.

Ich erhielt dieselbe am 11., wo Abends der Himmel hier ganz bedeckt war. Am 12. hingegen war es sehr heiter; ich fand den \textsc{Harding}'schen Stern sogleich und bestimmte seinen Ort durch eine gute Beobachtung. Seitdem habe ich ihn bis heute (den 30. Sept.) noch an eilf verschiedenen Abenden beobachten können, das letztemal am 28. Seine Bewegung ist noch immer rückläufig und nach Süden gehend, der bisherige Bogen kommt einer geraden Linie so nahe, dass die Krümmung kaum bemerkbar ist, die Geschwindigkeit hat aber merklich zugenommen. Diese Erscheinungen vertragen sich sehr gut mit der Voraussetzung, dass dieser Weltkörper wirklich eine planetarische Bahn beschreibe, berechtigen aber natürlich, so im allgemeinen genommen, nur erst zu Vermuthungen. Gewissheit kann uns nur ein
\SourceLeaf{260}
mit so vieler Schärfe und Vorsicht geführter Calcul geben, als zu einer so delicaten und
schlüpfrigen Untersuchung nöthig sind.

Das äussere Ansehen dieses Sterns ist dem der Ceres und Pallas ganz ähnlich, nur übertrifft
er sie noch etwas an Helligkeit. Nach meiner Schätzung kommt er einem Fixstern $7^{ter}$ bis
$8^{ter}$ Grösse ziemlich gleich, und ich habe ihn selbst am 18. September, wo der fast volle
Mond nicht weit davon entfernt war, mit einem \textsc{Short}ischen Spiegelteleskope ohne die
geringste Anstrengung sehen und beobachten können. Von einem Nebel, wie bei Cometen, zeigt sich
bei diesem Sterne, auch durch die grossen Lilienthaler Instrumente nicht die mindeste Spur.

Von andern auswärtigen Astronomen, die diesen Stern schon beobachten, sind mir bisher nur
Hr. Dr. \textsc{Olbers}, Hr. von \textsc{Zach} und Hr. Prof. \textsc{Bode} bekannt; der erste hat
die Beobachtungen am 7. Hr. von \textsc{Zach} am 13. und Hr. \textsc{Bode} am 21. angefangen. Da
indessen der Stern noch bis wenigstens zum Januar 1805 sichtbar bleiben wird, so habe ich die
zuversichtliche Hoffnung, dass er noch auf allen europäischen Sternwarten wird beobachtet werden
können. So kurze Zeit auch bisher dieser merkwürdige Weltkörper beachtet worden ist, so habe
ich doch bereits den Versuch gemacht, grösstentheils nach meinen eigenen Beobachtungen eine erste
genäherte Bestimmung von den Hauptdimensionen seiner wahren Laufbahn zu machen. Dies dient
theils zu einer vorläufigen Befriedigung der sehr natürlichen allgemeinen Wissbegierde; einen
wesentlichern Nutzen aber schafft eine solche Arbeit dadurch, dass sie uns in den Stand setzt, den
Lauf des Sternes schon auf einige Zeit \emph{voraus} zu sagen, und dadurch entferntern Astronomen,
welche die Nachrichten vom zurückgelegten Laufe erst so spät erhalten können, dass sie
unmittelbar zur Auffindung nicht mehr zulänglich sind. den Ort, wo sie ihn zu suchen haben,
beiläufig anzuzeigen, und so vielleicht fruchtloses Nachsuchen zu ersparen.

Das merkwürdige Resultat dieser Untersuchung ist gewesen, dass der \textsc{Harding}'sche Stern,
gleich der Ceres und Pallas, eine planetarische Bahn zwischen dem Mars und Jupiter beschreibt;
dass die Umlaufszeit und mittlere Entfernung von der Sonne fast (oder ganz) genau dieselben sind,
wie bei jenen; dass die Excentricität noch etwas grösser ist, als bei der Pallas, die Neigung
der Bahn hingegen viel kleiner, als bei dieser, und nur wenig grösser, als bei der Ceres. Es
ergab sich nemlich: Umlaufszeit 1662 Tage; mittlere Entfernung 2.746; Excentricität 0.287;
Neigung der Bahn 13 Grad 35 Minuten; Länge der Sonnenferne 239 Grad 14 Min.; Länge des
aufsteigenden Knoten 171 Grad 16 Min.; Mittlere Länge in der Bahn den 5. Sept. um Mittag im
Meridian der Sternwarte Seeberg 20 Grad 39 Min.

Diesen Resultaten zufolge, lässt sich nun schon als ausgemacht ansehen, dass der
\textsc{Harding}'sche Stern, gerade so, wie die andern Planeten, seine regelmässigen Perioden
der Sichtbarkeit und Unsichtbarkeit haben, und von jetzt an ein beständiger Gegenstand unsrer
Beobachtung sein wird. Er verdient daher den Planetenrang mit allem Rechte, und Herr
\textsc{Harding} hat deshalb um ihn den ältern Planeten gleich zu stellen, bereits den Namen
\emph{Juno} für ihn vorgeschlagen; es ist kein Zweifel, dass alle Astronomen das Recht des
Entdeckers ehren und seiner Benennung beitreten werden. Als Zeichen der Juno ist ihr Scepter der
einen Stern trägt in Vorschlag gebracht worden.

Die Juno steht jetzt innerhalb des sehr kenntlichen Vierecks, welches die vier Sterne in den
Fischen bilden, die in \textsc{Flamsteeds} Verzeichniss mit 27, 29, 30, 33 bezeichnet sind: gute
Augen sehen diese Sterne bei heiterm Himmel ohne Fernrohr. Die Juno wird ihre Bewegung noch einige
Zeit in der bisherigen Richtung, aber mit schon jetzt wieder abnehmender Geschwindigkeit
fortsetzen: in der letzten Hälfte des Octobers wird die Krümmung des Weges gegen Südosten
immer merklicher: etwa den 30. October verwandelt sich die rückläufige Bewegung in eine
rechtläufige, und etwa am 17. November wo die Juno eine südliche Declination von
$11\frac{1}{4}$ Grad erreicht haben wird, fängt diese wieder an abzunehmen, und die Richtung des
Weges krümmt sich wieder nordwärts.

\SourceLeaf{261}
Da vielleicht unter den Lesern dieser Blätter auch dieser und jener ist, der hinlängliche Werkzeuge und Geschicklichkeit besitzt, um sich mit Aufspürung dieses neuen Planeten zu versuchen, und ihn auf seinem Laufe zu verfolgen, so füge ich noch die genaue Angabe, sowol der bisherigen Stellungen als seiner künftigen bei, jene nur von Woche zu Woche, diese von 3 zu 3 Tagen. Bei allen ist die Mitternachtsstunde zu verstehen. [Diese Ephemeride ist ein Auszug aus der oben S. 253 abgedruckten.]
